% !TEX program = lualatex
\documentclass[professionalfonts]{beamer}
\newcommand{\slidemode}{1}  % カラー投影モード
\usepackage{beamer_template}
\usepackage{enumerate}

\newcommand{\LectureNum}{7}
\newcommand{\LectureTitle}{線形システムの伝達関数とデルタ関数}
\newcommand{\CourseName}{応用数学}
\renewcommand{\TermName}{前期}

\begin{document}

\begin{frame}{第7回　線形システムの伝達関数とデルタ関数}
  \begin{block}{今回の内容}
  線形システムの伝達関数とデルタ関数
  \end{block}
  \vfill\hfill{\scriptsize \textcolor{gray}{第2章 ラプラス変換　§2}}
\end{frame}

\begin{frame}{例題9}
  \begin{exampleblock}{問題}
    微分方程式 $x'' + ax' + bx = y(t)$ （ a, b は定数 $, x(0)=0, x'(0)=0$ ）に対して\\[2pt]
    線形システムの伝達関数 $H(s)$ を求めよ\\[2pt]
  \end{exampleblock}
\end{frame}

\begin{frame}{例題9　解答}
  初期値ゼロで両辺をラプラス変換:\\[4pt]
  $(s^{2}+as+b)X(s) = Y(s)$\\[4pt]
  $H(s) = X(s)/Y(s) = 1/(s^{2}+as+b)$
\end{frame}

\begin{frame}{演習問題}
  \begin{exampleblock}{自分で解いてみよう}
  \begin{description}
    \item[問10] 微分方程式 $x'' + ax' + bx = y(t), x(0)=0, x'(0)=0$ の
    伝達関数 $H(s)$ を求めよ\\
  \end{description}
  \end{exampleblock}
\end{frame}

\begin{frame}{問10　解答}
  初期値ゼロで両辺をラプラス変換:\\[4pt]
  $(s^{2}+as+b)X(s) = Y(s)$\\[4pt]
  $H(s) = X(s)/Y(s) = 1/(s^{2}+as+b)$
  \\[8pt]\textcolor{red}{\textbf{答}}：$H(s) = 1/(s^{2}+as+b)$
\end{frame}

\begin{frame}{演習問題}
  \begin{exampleblock}{自分で解いてみよう}
  \begin{description}
    \item[問11] $f(t)*\delta(t) = \delta(t)*f(t) = f(t)$ を、たたみこみの性質を用いて証明せよ
    \item[問12] 次の微分方程式の解を求めよ。ただし $\omega$ は正の定数とする。
    $y'' + \omega^{2}y = \delta(t),  y(0) = 0,  y'(0) = 0$\\
  \end{description}
  \end{exampleblock}
\end{frame}

\begin{frame}{問11　解答}
  ラプラス変換を使った証明:\\[4pt]
  $L[f(t)*\delta(t)] = F(s)\cdot L[\delta(t)] = F(s)\cdot1 = F(s) = L[f(t)]$\\[4pt]
  $\therefore f(t)*\delta(t) = f(t)$\\[4pt]
  交換法則より $\delta(t)*f(t) = f(t)*\delta(t) = f(t)$ （証明終）\\[4pt]
  （直接証明） $f(t)*\delta(t) = \int_{0}^t f(t-\tau)\delta(\tau)d\tau = f(t-0) = f(t)$
\end{frame}

\begin{frame}{問12　解答}
  両辺をラプラス変換（初期値ゼロ）:\\[4pt]
  $(s^{2} + \omega^{2})Y(s) = L[\delta(t)] = 1$\\[4pt]
  $Y(s) = 1/(s^{2}+\omega^{2}) = (1/\omega)\cdot\omega/(s^{2}+\omega^{2})$\\[4pt]
  逆ラプラス変換 $: y(t) = (1/\omega)\sin(\omega t)$
  \\[8pt]\textcolor{red}{\textbf{答}}：$y(t) = (1/\omega)\sin(\omega t)$（デルタ入力に対するインパルス応答）
\end{frame}

\end{document}
